\begin{helper}
Fix \(U_0,V_0\subseteq\operatorname{Fin}n\), \(t\), a blue graph \(G_B\), a
map \(\rho:\operatorname{Fin}n\to\operatorname{Fin}m\), and an exposure
\(\eta:U_0\to\operatorname{Fin}m\).  The conditional probability appearing in
\(\noderef{OppositeProjectionBlueExposureTailBound}\), where the blue graph
and second injection coordinate are conditioned and the first coordinate is
tested, is equal to the corresponding red-coordinate conditional probability
obtained by swapping the two graph coordinates and the two coordinates of
every injection value.
\end{helper}
\begin{proof}
Define an involution on two-bite samples by sending
\((G_R,G_B,\phi)\) to \((G_B,G_R,\phi^\sigma)\), where
\(\phi^\sigma(v)=((\phi(v))_2,(\phi(v))_1)\).  This map is a bijection of the
finite sample space.  Under it, the blue conditioning event becomes the red
conditioning event with \(G_R=G_B\), the second-coordinate tail event becomes
the first-coordinate tail event, and conversely.

The sample weight is invariant under this involution, because
\(\noderef{TwoBiteSampleWeight}\) is the product of the two graph weights and
the uniform injection weight, and multiplication is commutative.  Therefore
\(\noderef{TwoBiteEventProbability}\) is preserved for each event after
precomposition with the involution.  Applying this to the conditioning event
and to the intersection of the tail event with the conditioning event, and
then unfolding \(\noderef{TwoBiteConditionalProbability}\), gives equality of
the two conditional probabilities.
\end{proof}
